\documentclass[UTF8,11pt]{ctexart}
\usepackage[a4paper,margin=22mm]{geometry}
\usepackage{graphicx,enumitem,amssymb,tabularx}
\usepackage[normalem]{ulem}
\usepackage{xeCJKfntef}
\setlength{\parindent}{0pt}
\setlength{\parskip}{.6em}
\sloppy
\begin{document}
\section*{2024年6月英语四级真题(第2套)}
\subsection*{2024 年6 月大学英语四级考试真题（第2 套）}


\subsection*{\textbf{Part I Writing} \textbf{(30 minutes)}}


\textbf{Directions:} \textit{Suppose your university is seeking students’ opinions on whether university canteens should be open} \textit{to the public. You are now to write an essay to express your view. You will have 30 minutes for the task. You should} \textit{write at least }\textit{\uline{120 }}\textit{words but no more than }\textit{\uline{180 }}\textit{words.}


\subsection*{\textbf{Part II} \textbf{Listening Comprehension} \textbf{(25 minutes)}}


\subsection*{\textbf{Section A}}


\textbf{Directions:} \textit{In this section, you will hear three news reports. At the end of each news report, you will hear two or} \textit{three questions. Both the news report and the questions will be spoken only once. After you hear a question, you} \textit{must choose the best answer from the four choices marked A), B), C) and D). Then mark the corresponding letter} \textit{on }\textbf{\textit{Answer Sheet 1}}\textit{ with a single line through the centre.}


\textbf{Questions 1 and 2 are based on the news report you have just heard.}


1.


\begin{itemize}[leftmargin=2em,itemsep=.2em,topsep=.4em]\item[A.] Due to a fire alarm in their apartments.
\item[B.] Because of the smoke and heat damage.
\item[C.] Due to the water used to extinguish the flames.
\item[D.] Because of the collapse of the three-story building.
\end{itemize}

2.


\begin{itemize}[leftmargin=2em,itemsep=.2em,topsep=.4em]\item[A.] Investigating the cause of the incident.
\item[B.] Helping search for the suspect of the crime.
\item[C.] Rescuing the businessmen trapped in the building.
\item[D.] Checking town records for the property developer.
\end{itemize}

\textbf{Questions 3 and 4 are based on the news report you have just heard.}


3.


\begin{itemize}[leftmargin=2em,itemsep=.2em,topsep=.4em]\item[A.] It plays a less important role in one’s health than nutrient intake.
\item[B.] It impacts people’s health to a lesser degree than sun exposure.
\item[C.] It is associated with people’s mental health conditions.
\item[D.] It is linked with older adults’ symptoms of depression.
\end{itemize}

\noindent\begin{tabularx}{\linewidth}{@{}lXXXX@{}}4. & \textbf{A.} It was indefinite. & \textbf{B.} It was systematic. & \textbf{C.} It was straightforward. & \textbf{D.} It was insignificant.\\\end{tabularx}\par

\textbf{Questions 5 to 7 are based on the news report you have just heard.}


5.


\begin{itemize}[leftmargin=2em,itemsep=.2em,topsep=.4em]\item[A.] It has helped solve several murder cases.
\item[B.] It has become a star police dog in Beijing.
\item[C.] It has surpassed its mother in performance.
\item[D.] It has done better than naturally born dogs.
\end{itemize}

6.


\begin{itemize}[leftmargin=2em,itemsep=.2em,topsep=.4em]\item[A.] To speed up investigation into criminal cases.
\item[B.] To test the feasibility of cloning technology.
\item[C.] To cut down training expenses.
\item[D.] To reduce their training time.
\end{itemize}

7.


\begin{itemize}[leftmargin=2em,itemsep=.2em,topsep=.4em]\item[A.] Cloning is too complicated a process.
\item[B.] The technology is yet to be accepted.
\item[C.] Cloning is ethically controversial.
\item[D.] The technology is too expensive.
\end{itemize}

\subsection*{\textbf{Section B}}


\textbf{Directions:} \textit{In this section, you will hear two long conversations. At the end of each conversation, you will hear} \textit{four questions. Both the conversation and the questions will be spoken only once. After you hear a question, you} \textit{must choose the best answer from the four choices marked A), B), C) and D). Then mark the corresponding letter} \textit{on }\textbf{\textit{Answer Sheet 1}}\textit{ with a single line through the centre.}


\textbf{Questions 8 to 11 are based on the conversation you have just heard.}


8.


\begin{itemize}[leftmargin=2em,itemsep=.2em,topsep=.4em]\item[A.] He read it somewhere online.
\item[B.] He heard about it from a coworker.
\item[C.] He read an article reviewing it.
\item[D.] He watched a TV series based on it.
\end{itemize}

\noindent\begin{tabularx}{\linewidth}{@{}lXXXX@{}}9. & \textbf{A.} His publications. & \textbf{B.} His first book. & \textbf{C.} His address. & \textbf{D.} His name.\\\end{tabularx}\par

10.


\begin{itemize}[leftmargin=2em,itemsep=.2em,topsep=.4em]\item[A.] Collect a lot more data.
\item[B.] Relax a bit less often.
\item[C.] Clarify many new concepts.
\item[D.] Read more reference books.
\end{itemize}

11.


\begin{itemize}[leftmargin=2em,itemsep=.2em,topsep=.4em]\item[A.] Find out the show’s most interesting episodes.
\item[B.] Watch the series together with the woman.
\item[C.] Get an e-copy of the book to read.
\item[D.] Check to see when the show starts.
\end{itemize}

\textbf{Questions 12 to 15 are based on the conversation you have just heard.}


12.


\begin{itemize}[leftmargin=2em,itemsep=.2em,topsep=.4em]\item[A.] To check the prices of his farm produce.
\item[B.] To ask the way to the Newcastle City Hall.
\item[C.] To inquire about the vegetarian food festival.
\item[D.] To seek the man’s help with her work on the farm.
\end{itemize}

\noindent\begin{tabularx}{\linewidth}{@{}lXXXX@{}}13. & \textbf{A.} Bakers. & \textbf{B.} Vendors. & \textbf{C.} Vegetarians. & \textbf{D.} Organizers.\\\end{tabularx}\par

14.


\begin{itemize}[leftmargin=2em,itemsep=.2em,topsep=.4em]\item[A.] The issuing of certificates to vendors.
\item[B.] The completion of the baking task.
\item[C.] The festival they are organizing.
\item[D.] The deadline for application.
\end{itemize}

15.


\begin{itemize}[leftmargin=2em,itemsep=.2em,topsep=.4em]\item[A.] The closing date of submission
\item[B.] The website of his company.
\item[C.] The details of the ceremony.
\item[D.] The organizer’s address.
\end{itemize}

\subsection*{\textbf{Section C}}


\textbf{Directions:}\textit{ In this section, you will hear three passages. At the end of each passage, you will hear three or four} \textit{questions. Both the passage and the questions will be spoken only once. After you hear a question, you must} \textit{choose the best answer from the four choices marked A), B), C) and D). Then mark the corresponding letter on} \textbf{\textit{Answer Sheet 1}}\textit{ with a single line through the centre.}


\textbf{Questions 16 to 18 are based on the passage you have just heard.}


16.


\begin{itemize}[leftmargin=2em,itemsep=.2em,topsep=.4em]\item[A.] Most scenic sites have been closed.
\item[B.] Access to official campsites is limited.
\item[C.] Health experts advise going outdoors.
\item[D.] People have more time during the summer.
\end{itemize}

17.


\begin{itemize}[leftmargin=2em,itemsep=.2em,topsep=.4em]\item[A.] It is strongly opposed by nearby residents.
\item[B.] It leads to much waste of public money.
\item[C.] It has caused environmental concerns.
\item[D.] It has created conflicts among campers.
\end{itemize}

18.


\begin{itemize}[leftmargin=2em,itemsep=.2em,topsep=.4em]\item[A.] Look for open land in Scotland.
\item[B.] Leave no trace of their camping.
\item[C.] Avoid getting close to wilderness.
\item[D.] Ask for permission from authorities.
\end{itemize}

\textbf{Questions 19 to 21 are based on the passage you have just heard.}


19.


\begin{itemize}[leftmargin=2em,itemsep=.2em,topsep=.4em]\item[A.] They outcompete mythical creatures.
\item[B.] They usually mind their own business.
\item[C.] They truly exist in the Amazon region.
\item[D.] They resemble alarmingly large snakes.
\end{itemize}

20.


\begin{itemize}[leftmargin=2em,itemsep=.2em,topsep=.4em]\item[A.] Scar tissue from dolphins fighting.
\item[B.] Skin infection from water pollution.
\item[C.] Unhealed wounds from snake bites.
\item[D.] Swimming along in seasonal floods.
\end{itemize}

21.


\begin{itemize}[leftmargin=2em,itemsep=.2em,topsep=.4em]\item[A.] It has been shrinking at an astonishing pace.
\item[B.] It has been placed under international protection.
\item[C.] It has been appealing to both freshwater and sea dolphins.
\item[D.] It has been abandoned as a battleground for male dolphins.
\end{itemize}

\textbf{Questions 22 to 25 are based on the passage you have just heard.}


22.


\begin{itemize}[leftmargin=2em,itemsep=.2em,topsep=.4em]\item[A.] About 58\% of young adults call parental support the new normal.
\item[B.] Most adult children enjoy increasing sources of financial support.
\item[C.] A full 70\% of the young adults cannot afford to buy a car by themselves.
\item[D.] Most early adults cannot sustain their lifestyles without parental support.
\end{itemize}

23.


\begin{itemize}[leftmargin=2em,itemsep=.2em,topsep=.4em]\item[A.] It renders them dependent.
\item[B.] It causes them to lose dignity.
\item[C.] It makes them mentally immature.
\item[D.] It hinders them from getting ahead.
\end{itemize}

24.


\begin{itemize}[leftmargin=2em,itemsep=.2em,topsep=.4em]\item[A.] It challenges one’s willpower.
\item[B.] It results from education.
\item[C.] It calls for due assistance.
\item[D.] It defines adulthood.
\end{itemize}

\noindent\begin{tabularx}{\linewidth}{@{}lXXXX@{}}25. & \textbf{A.} Current lifestyles. & \textbf{B.} Poor budgeting & \textbf{C.} College loans. & \textbf{D.} Emergency expenses.\\\end{tabularx}\par

\subsection*{\textbf{Part III Reading Comprehension (40 minutes)}}


\subsection*{\textbf{Section A}}


\textbf{Directions:} \textit{In this section, there is a passage with ten blanks. You are required to select one word for each blank} \textit{from a list of choices given in a word bank following the passage. Read the passage through carefully before} \textit{making your choices. Each choice in the bank is identified by a letter. Please mark the corresponding letter for} \textit{each item on }\textbf{\textit{Answer Sheet 2}}\textit{ with a single line through the centre. You may not use any of the words in the bank} \textit{more than once.}


It’s well known that physical exercise is beneficial not just to physical health but also to mental health. Yet whereas most countries have \uline{26 }, evidence-backed guidelines on the type and intensity of exercise \uline{27} for various physical health benefits, such guidelines do not yet exist for exercise and mood. This is \uline{28} due to a lack of necessary evidence. However, a new systematic review brings us usefully up-to-date on the current findings in this area.


Before \uline{29} into some of the key take-aways, an important \uline{30} made in the review is between aerobic exercise and anaerobic. The former \uline{31} such things as walking, jogging and cycling and means exercising in such a way that your body is able to use oxygen to burn fat for energy. In contrast, anaerobic exercise — such as lifting heavy weights — is of such \uline{32} intensity that your body does not have time to use oxygen to create energy and so instead it breaks down glucose (葡萄糖) in your blood or muscles.


Beginning first with the influence of exercise intensity on the mood benefits of aerobic exercise, the researchers, led by John Chan at Shenzhen University, found \uline{33} results from 19 relevant studies. Some favoured higher intensity, others low, while seven studies found that intensity made no \uline{34} to mood benefits.


In relation to the intensity of anaerobic exercise, however, the results were far clearer — the optimum (最佳选择) for improving mood is \uline{35} intensity, perhaps because low intensity is too dull while high intensity is too unpleasant.


\begin{itemize}[leftmargin=2em,itemsep=.2em,topsep=.4em]\item[A.] constitutes
\item[B.] contradictory
\item[C.] decision
\item[D.] detailed
\item[E.] difference
\item[F.] dipping
\item[G.] distinction
\item[H.] falling
\item[I.] involves
\item[J.] moderate
\item[K.] notified
\item[L.] partly
\item[M.] required
\item[N.] traditionally
\item[O.] vigorous
\end{itemize}

\subsection*{\textbf{Section B}}


\textbf{Directions:} \textit{In this section, you are going to read a passage with ten statements attached to it. Each statement} \textit{contains information given in one of the paragraphs. Identify the paragraph from which the information is derived.} \textit{You may choose a paragraph more than once. Each paragraph is marked with a letter. Answer the questions by} \textit{marking the corresponding letter on }\textbf{\textit{Answer Sheet 2}}\textit{.} \textbf{Why Do Americans Work So Much?}


A) How will we all keep busy when we only have to work 15 hours a week? That was the question that worried the British economist John Maynard Keynes when he wrote his short essay “Economic Possibilities for Our Grandchildren” in 1930. Over the next century, he predicted, the economy would become so productive that people would barely need to work at all. For a while, it looked like Keynes was right. In 1930 the average working week was 47 hours in the United States. But by 1970, the number of hours Americans worked on average had fallen to slightly less than 39.


B) But then something changed. Instead of continuing to decline, the duration of the working week remained stable. It has stayed at just below 40 hours for nearly five decades. So what happened? Why are people working just as much today as in 1970?


C) There would be no mystery in this if Keynes had been wrong about the power of technology to increase the economy’s productivity, which he thought would lead to a standard of living “between four and eight times as high as it is today.” But Keynes got that right: Technology has made the economy massively more productive. According to Benjamin M. Friedman, an economist at Harvard, the U.S. economy is right on track to reach Keynes’s eight-fold (八倍) multiple by 2029. That is a century after the last data Keynes would have had access to.


D) In a new paper, Friedman tries to figure out why that increased productivity has not translated into increased leisure time. Perhaps people just never feel materially satisfied, always wanting more money to buy the next new thing. This is a theory that appeals to many economists. “This argument is, at best, far from sufficient,” he writes. If that were the case, why did the duration of the working week decline in the first place?


E) Another theory Friedman considers is that, in an era of ever fewer settings that provide effective opportunities for personal connections and relationships, people may place more value on the socializing that happens at work. There is support for this theory. Many people today consider colleagues as friends. But Friedman argues that the evidence for this theory is far from conclusive. Many workers report that they would like to spend more time with family, rather than at work. Furthermore, this theory cannot explain the change in trend in the U.S. working week in the 1970s.


F) A third possibility proves more convincing for Friedman. That is: American inequality means that the gains of increasing productivity are not widely shared by everyone. In other words, most Americans are too poor to work less. Unlike the other two explanations Friedman considers, this one fits chronologically (按年代). Inequality declined in America during the period following World War II, along with the duration of the working week. But since the early 1970s it has risen dramatically.


G) Keynes’s prediction of a shorter working week rests on the idea that the standard of living would continue rising for everyone. But Friedman says that this is not what has happened. Although Keynes’s eight-fold figure holds up for the economy as a whole, it is not at all the case for the median (中位数的) American worker. For them, output by 2029 is likely to be around 3.5 times what it was when Keynes was writing. This is a bit below his four- to eight-fold predicted range.


H) This can be seen in the median worker’s income over this time period, complete with a shift in 1973 that fits in precisely with when the working week stopped shrinking. According to Friedman, between 1947 and 1973 the average hourly wage for normal workers (those who were not in management roles) in private industries other than agriculture nearly doubled in terms of what their money could buy. But by 2013 the average hourly wage for ordinary workers had fallen 5 percent from the 1973 level in terms of actual purchasing power. Thus, though American incomes may have gone up since 1973, the amount that American workers can actually buy with their money has gone down. For most Americans, then, the magic of increasing productivity stopped working around 1973. Thus, they had to keep working just as much in order to maintain their standard of living.


I) What Keynes predicted was a very optimistic version of what economists call technological unemployment. This is the idea that less labor will be necessary because machines can do so much. In Keynes’s vision, the resulting unemployment would be distributed more or less evenly across society in the form of increased leisure. But Friedman says that, for Americans, reality is much darker. Americans now have a labor market in which millions of people — those with fewer skills and less education — are seeking whatever poorly paid work they can get. This is confirmed by a recent poll that found that, for half of hourly workers, their top concern is not that they work too much but that they work too little. This is most likely not because they like their jobs so much. Rather, we can assume it is because they need the money.


J) This explanation leaves an important question. If the very rich — the workers who have reaped above-average gains from the increased productivity since Keynes’s time — can afford to work less, why do they continue to work so much? (Indeed, research has shown that the highest earners in America tend to work the most.) Friedman believes that for many top earners, work is a labor of love. They are doing work they care about and are interested in, and doing more of it is not necessarily a burden. For them, it may even be a pleasure. These top earners derive meaning from their jobs and work is an important part of how they think of themselves. And, of course, they are compensated for it at a level that makes it worth their while.


K) Friedman concludes that the prosperity (繁荣) Keynes predicted is here. After all, the economy as a whole has grown even more brilliantly than he expected. But for most Americans, that prosperity is nowhere to be seen. And, as a result, neither are those shorter working weeks.


36. Some people view socializing at the workplace as a chance to develop personal relationships.


37. As ordinary American workers’ average hourly pay had decreased despite increasing productivity, they


had to work just as many hours as before to keep their living standards.


38. American workers’ average weekly working time has not changed for nearly half a century.


39. Friedman believes inequality in the U.S. largely explains why increasing productivity has not resulted in


reduced working hours.


40. Many economists assume people’s thirst for material things has prevented them from enjoying more


leisure time.


41. An economist’s prediction about a shorter average working week seemed to be correct for a time in the


20th century.


42. In the U.S. labor market, the primary concern of people with less schooling and fewer skills is to secure


any employment even if it is low-paid.


43. Keynes was right in predicting that technology would make the economy much more productive.


44. Many of the highest earners have a keen interest in and love for what they are doing.


45. According to Keynes, there would be a shorter working week with everyone’s standard of living


continuing to rise.


\subsection*{\textbf{Section C}}


\textbf{Directions: }\textit{There are 2 passages in this section. Each passage is followed by some questions or unfinished} \textit{statements. For each of them there are four choices marked A), B), C) and D). You should decide on the best} \textit{choice and mark the corresponding letter on }\textbf{\textit{Answer Sheet 2}}\textit{ with a single line through the centre.}


\subsection*{\textbf{Passage One}}


\textbf{Questions 46 to 50 are based on the following passage.}


Lao Zi once said, “Care about what other people think and you will always be their prisoner.”


People-pleasing, or seeking self-worth through others’ approval, is unproductive and an exhausting way to go through life. Why do we allow what others think of us to have so much power over how we feel about ourselves? If it’s true that you can’t please all people all of the time, wouldn’t it make sense to stop trying?


Unfortunately, sense often isn’t driving our behavior. For social beings who desire love and belonging, wanting to be liked, and caring about the effect we have on others, is healthy and allows us to make connections. However, where we get into trouble is when our self-worth is dependent upon whether we win someone’s approval or not.


This need to be liked can be traced back to when we were children and were completely dependent on others to take care of us: Small children are not just learning how to walk and communicate, they are also trying to learn how the world works. We learn about who we are and what is expected of us based on interactions with others, so, to a four-year-old, if Mommy or Daddy doesn’t like him or her, there is the danger that they will abandon them. We need to understand that when we desperately want someone to approve of us, it’s being driven by that little kid part of us that is still terrified of abandonment.


As you become more capable of providing yourself with the approval you seek, your need for external validation will start to vanish, leaving you stronger, more confident, and yes, happier in your life. Imagine how much time we lose each moment we restrain our authentic selves in an effort to be liked.


If we base our worth on the opinions of others, we cheat ourselves of the power to shape our experiences and embrace life not only for others but also for ourselves, because ultimately, there is no difference. So embrace the cliché (老话) and love yourself as it’s highly doubtful that you’ll regret it.


46. What can we conclude from Lao Zi’s quotation?


\begin{itemize}[leftmargin=2em,itemsep=.2em,topsep=.4em]\item[A.] We should see through other people’s attempt to make a prisoner of us.
\item[B.] We can never really please other people even if we try as hard as we can.
\item[C.] We can never be truly free if taking to heart others’ opinion of us.
\item[D.] We should care about other people’s view as much as they care about our own.
\end{itemize}

47. What will happen if we base our self-worth on other people’s approval?


\begin{itemize}[leftmargin=2em,itemsep=.2em,topsep=.4em]\item[A.] Our desire to be loved will be fulfilled.
\item[B.] Our life will be unfruitful and exhausting.
\item[C.] Our identity as social beings will be affected.
\item[D.] Our sense of self will be sharpened and enhanced.
\end{itemize}

48. What may account for our need to be liked or approved of?


\begin{itemize}[leftmargin=2em,itemsep=.2em,topsep=.4em]\item[A.] Our desperate longing for interactions with others.
\item[B.] Our understanding of the workings of the world.
\item[C.] Our knowledge about the pain of abandonment.
\item[D.] Our early childhood fear of being deserted.
\end{itemize}

49. What can we do when we become better able to provide ourselves with the desired approval?


\begin{itemize}[leftmargin=2em,itemsep=.2em,topsep=.4em]\item[A.] Enjoy a happier life.
\item[B.] Exercise self-restraint.
\item[C.] Receive more external validation.
\item[D.] Strengthen our power of imagination.
\end{itemize}

50. What does the author advise us to do in the last paragraph?


\begin{itemize}[leftmargin=2em,itemsep=.2em,topsep=.4em]\item[A.] Embrace life for ourselves and for others.
\item[B.] Base our worth on others’ opinions.
\item[C.] See our experiences as assets.
\item[D.] Love ourselves as we are.
\end{itemize}

\subsection*{\textbf{Passage Two}}


\textbf{Questions 51 to 55 are based on the following passage.}


Some people have said aging is more a slide into forgetfulness than a journey towards wisdom. However, a growing body of research suggests that late-in-life learning is possible. In reality, education does an aging brain good.


Throughout life, people’s brains constantly renovate themselves. In the late 1960s, British brain scientist Geoffrey Raisman spied growth in damaged brain regions of rats through an electron microscope; their brains were forging new connections. This meant brains may change every time a person learns something new.


Of course, that doesn’t mean the brain isn’t affected by the effects of time. Just as height usually declines over the years, so does brain volume: Humans lose about 4 percent every decade starting in their 40s. But that reduction doesn’t necessarily make people think slower; as long as we are alive and functioning, we can alter our brains with new information and experiences.


In fact, scientists now suspect accumulating novel experiences, facts, and skills can keep people’s minds more flexible. New pathways can strengthen our ever-changing mental structure, even as the brain shrinks.


Conventional fixes like word puzzles and brain-training apps can contribute to mental durability. Even something as simple as taking a different route to the grocery store or going somewhere new on vacation can keep the brain healthy.


A desire for new life challenges can further boost brainpower. Research about aging adults who take on new enterprises shows improved function and memory as well as a reduced risk of mental disease. Openness — a characteristic defined by curiosity and a desire for knowledge — may also help folks pass brain tests. Some folks are born with this take-in-the-world attitude, but those who aren’t as genetically gifted aren’t necessarily out of luck. While genes can encourage an interest in doing new things, a 2012 study in the journal \textit{Psychology and Aging} found completing reasoning tasks like puzzles and number games can enhance that desire for novel experiences, which can, in turn, refresh the brain. That’s why brain scientist Richard Kennedy says “It’s not that old dogs can’t learn new tricks. It’s that maybe old dogs don’t realize why they should.”


51. What do some people think of aging adults?


\begin{itemize}[leftmargin=2em,itemsep=.2em,topsep=.4em]\item[A.] Their wisdom grows as time goes by.
\item[B.] Their memory gradually deteriorates.
\item[C.] They can benefit from late-in-life learning.
\item[D.] They are likely to have mental health issues.
\end{itemize}

52. What can we conclude from Geoffrey Raisman’s finding?


\begin{itemize}[leftmargin=2em,itemsep=.2em,topsep=.4em]\item[A.] Brain damage seriously hinders one’s learning.
\item[B.] Brain power weakens slower than we imagine.
\item[C.] Brains can refresh and improve with learning.
\item[D.] Brains forge connections under new conditions.
\end{itemize}

53. What is one thing that helps maintain the health of our brain even as it shrinks?


\begin{itemize}[leftmargin=2em,itemsep=.2em,topsep=.4em]\item[A.] Doing daily routines by conventional means.
\item[B.] Avoiding worrying about our mental durability.
\item[C.] Imitating old dogs’ way of learning new tricks.
\item[D.] Approaching everyday tasks in novel ways.
\end{itemize}

54. What does the author say can contribute to the improvement of brain function?


\begin{itemize}[leftmargin=2em,itemsep=.2em,topsep=.4em]\item[A.] Being curious and desiring knowledge.
\item[B.] Being eager to pass brain tests at an old age.
\item[C.] Rising to life’s challenges and avoiding risks.
\item[D.] Boosting immunity to serious mental diseases.
\end{itemize}

55. What is the finding of the 2012 study in the journal \textit{Psychology and Aging}?


\begin{itemize}[leftmargin=2em,itemsep=.2em,topsep=.4em]\item[A.] Wishing to solve puzzles enhances one’s reasoning power.
\item[B.] Playing number games unexpectedly stimulates one’s memory.
\item[C.] Desiring new experiences can help to renovate the brain.
\item[D.] Learning new tricks should not be confined to old dogs only.
\end{itemize}

\subsection*{\textbf{Part IV Translation (30 minutes)}}


\textbf{Directions: }\textit{For this part, you are allowed 30 minutes to translate a passage from Chinese into English. You} \textit{should write your answer on }\textbf{\textit{Answer Sheet 2}}\textit{.}


\CJKunderline{农历 }(the lunar calendar) 起源于数千年前的中国，根据太阳和月亮的运行规律制定。长期以来，农历在农业生产和人们日常生活中发挥着重要作用。古人依据农历记录日期、安排农活，以便最有效地利用自然资源和气候条件，提高农作物的产量和质量。中国的春节、中秋节等传统节日的日期都基于农历。农历是中国传统文化的重要组成部分，当今依然广为使用。

\end{document}